\scanpage{233}
\mspage{179}
% archive sheet number: 232
Les propriétés des foncteurs $\Phi,\Psi,C$ qui ont été
utilisées, pour \deleted{\unsure{obtenir les}}\inserted{fournir} tous les quasi-inverses
voulus des foncteurs (entre cat. localisées) induits par les foncteurs d'inclusion du diagramme, sont les suivantes.

\textcircled{a} On a \unsure{un} lien fonctoriel dans $\mathrm{Cat}_{/S}$
\[
X\xrightarrow{i_X}\Phi(X),\qquad \underline{\underline{\unsure{j|i}_X\in W_S^{\unsure{d|1}}}} .
\]
\deleted{\unsure{qui est dans} $W_S^d$,} Quand $X$ est dans $\underline{\mathrm{Fib}}_W\,S$,
$\Phi(X)$ est dans $\underline{\mathrm{Fibsc}}_0\,\deleted{\unsure{S}}S\deleted{\unsure{S}}$,
\deleted{\unsure{(sa restriction), donc}}\inserted{donc $i_X\in W_S^f$ ($=W_S^d$ dans $\underline{\mathrm{Fib}}_W\,S$)}.
\inserted{Ainsi} $\Phi$ induit le foncteur $\Phi_0$ du diagramme.

Dualement, on a
\[
X\mathrel{\overset{i_X}{\unsure{\longrightarrow|\text{---}}}}\Psi(X),
\qquad\underline{i_X\in W_S^g},
\]
et quand $X$ est dans $\underline{\mathrm{Fib}}_W\,S$, $\Psi(X)$ est dans
$\underline{\mathrm{Cofibsc}}_0\,S$, d'où le foncteur $\Psi_0$.

\textcircled{b} On a un lien fonctoriel
\[
\left|\begin{aligned}
&X\mathrel{\overset{i_X}{\unsure{\text{---}|\longrightarrow}}}C(X),\qquad\underline{i_X\in W_S}.\\
&\text{et }C(X)\in\mathrm{Ob}\,\underline{\mathrm{Parf}}_W\,S .
\end{aligned}\right.
\]
\deleted{De plus, quand $X$ est dans $\underline{\mathrm{Fib}}_W\,S$, on a
\[
i_X\in W_S^f\quad\text{si}\quad X\in\mathrm{Ob}\,\underline{\mathrm{Fib}}_W\,S .
\]}
Mais on fera attention qu'en l'absence de l'axiome
\margin{$\vert\vert\vert$}$W(4\mathrm{b})$ sur $W$, on ne peut \uline{pas} affirmer $i_X\in W_S^f$,
même si $X$ est dans $\underline{\mathrm{Fib}}_W\,S$\deleted{\unsure{.}} (voire, dans $\underline{\mathrm{Parf}}_W\,S$),
seul cas (de toutes façons) où on est en droit de s'y
attendre. Je note à ce propos le diagramme \uline{cartésien},
pour $f:X\mathrel{\unsure{\text{---}|\longrightarrow}}\deleted{\unsure{C}}Y$ dans $\mathrm{Cat}_{/S}$
\diagram{\begin{tikzcd}[ampersand replacement=\&,column sep=large]
C(Y)_{\unsure{s|S}}\arrow[d,"\pi_Y^s"'] \& C(X)_s\arrow[l,"C(f)_s"']\arrow[d,"\pi_X^s"]\\
Y \& X\arrow[l,"f"']\;{,}
\end{tikzcd}}
\margin{flèches verticales dans $\mathrm{Parf}_W\subset\mathrm{Fib}_W$}
qui montre par contre que \deleted{si $i_X,i_Y\in W_S^f$, \unsure{alors} $f\in W_S$,}
